\documentclass[10pt]{article}
\usepackage[T1]{fontenc}
\usepackage[utf8]{inputenc}
\usepackage[french]{babel}
\usepackage{amsmath,amssymb,mathrsfs,mathpazo}
\usepackage{geometry}
\geometry{a4paper,margin=2cm}
\usepackage{graphicx,array,tabularx,tabularray,colortbl,multirow,multicol,booktabs,enumitem,pifont,ulem,fancybox,diagbox,cancel}
\usepackage{pstricks,pstricks-add,pst-plot,pst-tree,pst-node,pst-text,pst-func,pst-math,pst-bspline,pst-3dplot,pst-eucl,pst-solides3d}
\let\pscalculate\undefined
\usepackage{pst-bezier}
\usepackage{tikz,pgfplots,tkz-tab}
\usetikzlibrary{arrows,arrows.meta,patterns,calc,positioning,trees,decorations.markings,shapes.geometric,shapes.misc,shapes.arrows,intersections,decorations.pathmorphing,decorations.pathreplacing,backgrounds,shadows,babel}
\pgfplotsset{compat=1.18}
\usepgfplotslibrary{fillbetween}
\usepackage[np]{numprint}
\usepackage{listings,makecell,pas-tableur,twoopt}
\let\barre\undefined
\usepackage[pstricks]{tablvar}
\definecolor{darkgreen}{RGB}{0,100,0}
\providecommand{\red}{\color{red}}
\providecommand{\blue}{\color{blue}}
\providecommand{\green}{\color{green}}
\providecommand{\black}{\color{black}}
\providecommand{\white}{\color{white}}
\providecommand{\gray}{\color{gray}}
\providecommand{\cyan}{\color{cyan}}
\providecommand{\magenta}{\color{magenta}}
\providecommand{\yellow}{\color{yellow}}
\providecommand{\violet}{\color{violet}}
\providecommand{\orange}{\color{orange}}
\providecommand{\euro}{\texteuro}
\providecommand{\EUR}{\texteuro}
\providecommand{\ieme}{\textsuperscript{e}}
\providecommand{\degres}{\ensuremath{^\circ}}
\providecommand{\degree}{\ensuremath{^\circ}}
\providecommand{\vv}[1]{\overrightarrow{#1}}
\providecommand{\Coord}[1]{#1}
\providecommand{\decosix}{}
\providecommand{\starredbullet}{\ensuremath{\bullet}}
\pagestyle{empty}
\setlength{\parindent}{0pt}
\usepackage{adjustbox}\DeclareUnicodeCharacter{2236}{:}\DeclareUnicodeCharacter{2190}{\ensuremath{\leftarrow}}
\providecommand{\R}{\mathbb{R}}
\providecommand{\N}{\mathbb{N}}
\providecommand{\Z}{\mathbb{Z}}
\providecommand{\Q}{\mathbb{Q}}
\providecommand{\D}{\mathbb{D}}
\providecommand{\C}{\mathbb{C}}
\providecommand{\e}{{\mathrm{e}}}
\providecommand{\vect}[1]{\overrightarrow{#1}}
\providecommand{\vectt}[1]{\overrightarrow{\mathrm{#1}}}
\providecommand{\vvt}[1]{\overrightarrow{\mathrm{#1}}}
\providecommand{\np}[2][]{{#2~#1}}
\providecommand{\vv}[1]{{\overrightarrow{#1}}}
\renewcommand{\d}{{\mathrm{d}}}
\providecommand{\degres}{{^\circ}}
\providecommand{\degre}{{^\circ}}
\providecommand{\blue}{\color{blue}}
\providecommand{\red}{\color{red}}
\providecommand{\enskip}{\quad}
\providecommand{\llbracket}{\lbrack\!\lbrack}
\providecommand{\rrbracket}{\rbrack\!\rbrack}
\providecommand{\euro}{{\text{€}}}
\providecommand{\no}{{\text{n°}}}
\providecommand{\interval}[2]{{[#1;#2]}}
\providecommand{\intervFO}[2]{{[#1;#2[}}
\providecommand{\dotfill}{\ldots}

\providecommand{\ds}{\displaystyle}
\providecommand{\textding}[1]{#1}
\providecommand{\Oij}{\(\mathrm{O};\vec{\imath},\vec{\jmath}\)}
\providecommand{\Oijk}{\(\mathrm{O};\vec{\imath},\vec{\jmath},\vec{k}\)}
\providecommand{\Ouv}{\(\mathrm{O};\vec{u},\vec{v}\)}
\providecommand{\intervFF}[2]{\ensuremath{[#1;#2]}}
\renewcommand{\euro}{\ifmmode\text{\texteuro}\else\texteuro\fi}
\DeclareUnicodeCharacter{00B0}{\ifmmode^{\circ}\else\textdegree\fi}
\DeclareUnicodeCharacter{20AC}{\ifmmode\text{\texteuro}\else\texteuro\fi}
% Énoncés : Éducation nationale. Transcription et corrigés : APMEP, Denis Vergès et contributeurs. Découpage et classement : CoopMaths, Sébastien Lozano, Éric Elter et contributeurs. Adaptation de la présentation : Maths Explorer.
% CC BY-SA : https://coopmaths.fr/www/about/
% https://coopmaths.fr/alea/static/eam/2026/tex/eam_specialite_2026_06_asie_2_cor.tex
\begin{document}
\noindent\textit{APMEP / CoopMaths. CC BY-SA. Présentation adaptée par Maths Explorer.}\par\medskip
\subsection*{Partie A -- Étude d'un cercle}
	\begin{enumerate}
		\item Un point {M~$(x;y)$} appartient à $\mathcal{C}$ si et seulement si la distance AM (donc la norme du vecteur $\vectt{AM}$) est égale à {$\dfrac{\sqrt{37}}{2}$.}

		Il faut donc que {$\left\|\overrightarrow{\text{AM}}\right\|^2=\left(\dfrac{\sqrt{37}}{2}\right)^2=\dfrac{37}{4}$}

		Les coordonnées de {$\vectt{AM}$} sont {$\left(\begin{array}{c}x-3\\ y-1\end{array}\right)$} et donc {$\left\|\overrightarrow{\text{AM}}\right\|^2=(x-3)^2+(y-1)^2$.}

		Une équation cartésienne du cercle {$\mathcal{C}$} est donc {$(x-3)^2+(y-1)^2=\dfrac{37}{4}$.}

		\item On vérifie si les coordonnées de B vérifient cette équation :\quad  {$(6-3)^2+(\np{0.5}-1)^2=3^2+(-\np{0.5})^2=9+\np{0.25}=\np{9.25}$}.

		Or \quad {$\np{9.25}=\dfrac{37}{4}$}, \quad donc B appartient à {$\mathcal{C}$.}
	\end{enumerate}

\subsection*{Partie B -- Étude d'un cas particulier}

	\begin{enumerate}
		\item L'équation de {$\mathcal{H}$} est {$y=\dfrac 3 x$}, donc pour {$x=6$}, on obtient {$y=\dfrac 3 6=\np{0.5}$}, qui est bien l'ordonnée de B, qui appartient donc à {$\mathcal{H}$.}

		\item \begin{enumerate}
			\item Les coordonnées de {$\vectt{AD}$} sont {$\left(\begin{array}{c}1\\ 1\end{array}\right)$}, et celles de {$\vectt{AB}$} sont {$\left(\begin{array}{c}3\\ -\np{0.5}\end{array}\right)$.}

			\item On a donc {$\vectt{AD}\cdot\vectt{AB}=1\times3+1\times(-\np{0.5})=\np{2.5}$.}

			\item L'angle  $\widehat{\text{DAB}}$ n'est pas un angle droit, en effet s'il l'était, le produit scalaire {$\vectt{AD}\cdot\vectt{AB}$} serait nul.
		\end{enumerate}
	\end{enumerate}

\subsection*{Partie C -- Étude du cas général}

	\begin{enumerate}
		\item On calcule \quad {$\Delta=(-4)^2-4\times1\times3=4=2^2$}.

		Le discriminant est strictement positif, l'équation a deux solutions dans $\R$ qui sont {$\dfrac{4-2}{2}=1$} et {$\dfrac{4+2}{2}=3$}.

		\item $\widehat{\text{DAM}}$ est un angle droit si et seulement si A et M sont deux points distincts vérifiant {$\vectt{AD}\cdot\vectt{AM}=0$}.

		(les points doivent être distincts, sinon A et M ne définissent pas une droite, et donc l'angle $\widehat{\text{DAM}}$ n'est pas défini).

		On connaît déjà les coordonnées de {$\vectt{AD}$}, et celles de {$\vectt{AM}$} sont {$\left(\begin{array}{c}x-3\\ y-1\end{array}\right)$.}

		La condition s'écrit donc \quad {$x-3+y-1=0$}.

		Sachant que \quad {$y=\dfrac3 x$},\quad cette équation devient \quad {$x-3+\dfrac 3 x -1=0$} \quad soit, après multiplication par $x\neq0$ :
		\quad
		{$x^2-4x+3=0$}

		On sait que cette équation a deux solutions : la solution $x=3$ est évidemment sans intérêt : le produit scalaire dans ce cas est nul puisque A et M sont confondus.

		Il reste la solution $x=1$ : le point {M~$(1;3)$} est le seul point de {$\mathcal{H}$} tel que $\widehat{\text{DAM}}$ soit un angle droit.
	\end{enumerate}

\end{document}
